\subsection{Tasks and Datasets} 
We conduct evaluations of our \model and diverse baselines on a range of downstream tasks, holistically evaluating outputs with metrics designed to assess overall correctness, factuality, and fluency. 
Throughout these experiments, we conduct zero-shot evaluations, where we provide instructions describing tasks without few-shot demonstrations~\citep{wei2022finetuned,sanh2022multitask}. 
Details of our experiments' settings, including test-time instructions, are available in the Appendix Section~\ref{sec:training_details}. 


\noindent{\bf Closed-set tasks} include two datasets, i.e., a fact \textit{verification dataset} about public health ({\bf PubHealth}; \citealt{zhang2023interpretable}) and a \textit{multiple-choice reasoning dataset} created from scientific exams ({\bf ARC-Challenge}; ~\citealt{clark2018think}). 
We use accuracy as an evaluation metric and report on the test set. 
We aggregate the answer probabilities of target classes for both of these datasets (Appendix Section~\ref{sec:details_experiments}). 
 

\noindent{\bf Short-form generations tasks} include two open-domain question answering (QA) datasets, PopQA~\citep{mallen2022not} and TriviaQA-unfiltered~\citep{joshi-etal-2017-triviaqa}, where systems need to answer arbitrary questions about factual knowledge. 
For PopQA, we use the long-tail subset, consisting of 1,399 rare entity queries whose monthly Wikipedia page views are less than 100. 
As the TriviaQA-unfiltered (open) test set is not publicly available, we follow prior work's validation and test split~\citep{min-etal-2019-discrete,guu2020retrieval}, using 11,313 test queries for evaluation. 
We evaluate performance based on whether gold answers are included in the model generations instead of strictly requiring exact matching, following \cite{mallen2022not,schick2023toolformer}. 

\noindent{\bf Long-form generation tasks} include a biography generation task~\citep{min2023factscore} and a long-form QA task~{\bf ALCE-ASQA}~\cite{gao2023enabling,stelmakh2022asqa}. 
We use FactScore~\citep{min2023factscore} to evaluate biographies, and we use official metrics of correctness (str-em), fluency based on MAUVE~\citep{pillutla2021MAUVE}, and citation precision and recall~\citep{gao2023enabling} for ASQA. \footnote{\url{https://github.com/princeton-nlp/ALCE}} 

\subsection{Baselines}
\noindent{\bf Baselines without retrievals.}
We evaluate strong publicly available pre-trained LLMs, Llama2$_{\textsc{7b},\textsc{13b}}$~\citep{touvron2023llama}, instruction-tuned models, Alpaca$_{\textsc{7b},\textsc{13b}}$~\citep{dubois2023alpacafarm} (our replication based on Llama2); and models trained and reinforced using private data, ChatGPT~\citep{ouyang2022training} and Llama2-chat$_{\textsc{13b}}$. 
For instruction-tuned LMs, we use the official system prompt or instruction format used during training if publicly available. 
We also compare our method to concurrent work, CoVE$_\textsc{65b}$~\citep{dhuliawala2023chainofverification}, which introduces iterative prompt engineering to improve the factuality of LLM generations. 

\noindent{\bf Baselines with retrievals.}  
We evaluate models augmented with retrieval at test time or during training. The first category includes standard RAG baselines, where an LM (Llama2, Alpaca) generates output given the query prepended with the top retrieved documents using the same retriever as in our system. It also includes {Llama2-FT}, where Llama2 is fine-tuned on all training data we use without the reflection tokens or retrieved passages.
{We also report the result of retrieval-augmented baselines with LMs trained with private data: Ret-ChatGPT and Ret-Llama2-chat, which deploy the same augmentation technique above, as well as perplexity.ai, an InstructGPT-based production search system. 
}
The second category includes concurrent methods that are trained with retrieved text passages, i.e., { SAIL}~\citep{luo2023sail} to instruction-tune an LM on the Alpaca instruction-tuning data with top retrieved documents inserted before instructions, and {Toolformer}~\citep{schick2023toolformer} to pre-train an LM with API calls (e.g., Wikipedia APIs).\footnote{We report numbers using the results reported in the paper as the implementations are not available.}


\subsection{Experimental settings}
\noindent{\bf Training data and settings.}
Our training data consists of diverse instruction-following input-output pairs. 
In particular, we sample instances from Open-Instruct processed data~\citep{wang2023far} and knowledge-intensive datasets~\citep{petroni-etal-2021-kilt,stelmakh2022asqa,mihaylov-etal-2018-suit}. 
In total, we use 150k instruction-output pairs. 
We use Llama2 7B and 13B~\citep{touvron2023llama} as our generator base LM, and we use Llama2 7B as our base critic LM. 
For the retriever model \mret, we use off-the-shelf Contriever-MS MARCO~\citep{izacard2021towards} by default and retrieve up to ten documents for each input. 
More training details are in the Appendix Section~\ref{sec:training_details}. 

\noindent{\bf Inference settings.}
As a default configuration, we assign the weight terms \crel, \cgr, \cuse values of 1.0, 1.0 and 0.5, respectively.
 To encourage frequent retrieval, we set the retrieval threshold to 0.2 for most tasks and to 0 for ALCE~\citep{gao2023enabling} due to citation requirements.    
% in the subsequent analysis section, we examine the impacts of these varying hyperparameters.
We speed up inference using vllm~\citep{kwon2023efficient}. 
At each segment level, we adopt a beam width of 2.  
For a token-level generation, we use greedy decoding.  
By default, we use the top five documents from Contriever-MS MARCO~\citep{izacard2021towards}; for biographies and open-domain QA, we use additional top five documents retrieved by a web search engine, following~\cite{luo2023sail}; for ASQA, we use the author-provided top 5 documents by GTR-XXL~\citep{ni2021large} across all baselines for a fair comparison. 



